% Pillar-5 (MIN-REPORT-RL) conclusion + call to action
\section{Conclusion: A Call to Action for Venues}
\label{sec:p5-conclusion}
The evidence pattern is consistent and, we argue, sufficient to change reporting
policy: on a fixed stack, four algorithm labels spanned $0.033$ reward (not
shown equivalent at 3 seeds); across under-specified stacks, a nominal backend
comparison moved reward by $0.806$, of which $0.519$ was a base-checkpoint
switch and the rest a multiply-confounded harness difference; and one label
was attached to a mechanism it does not name. The backend gap is not a causal
estimate; both are evidence that RL-for-LLM results
are stack-conditioned. The community's papers should say so, in a form
machines can check.

Our call to action is addressed to venues, because venues set the price of
non-disclosure:
\begin{enumerate}[leftmargin=1.5em, itemsep=2pt, topsep=2pt]
\item \textbf{Require the manifest.} Any paper whose contribution depends on
an RL-for-LLM comparison should attach the seven-field run manifest and the
eighth evaluation report (pass@$k$; \tableref{tab:p5-minreport}) per compared
run as supplementary JSON.
Auto-emission (\secref{sec:p5-toolchain}) is a proposed integration, not a
demonstrated one-line cost. For closed stacks, ``provider-internal, version-dated'' is an
acceptable and informative value---the point is that conditioning be
declared, not that all stacks be open.
\item \textbf{Grade the comparisons.} Reviewers should see the
\texttt{grpo-stackdiff} grade and coverage for every headline comparison.
A comparison whose only differences are the defining R4 fields plus R1
nuisance identifies fewer known confounds but does not alone support an
algorithmic claim; experimental design must also be checked. Any additional
R2, R3, or R5 difference should be rewritten as a stack-conditioned
observation, and an underspecified pair should be treated as a
missing-experiment flag, exactly as a missing baseline is today.
\item \textbf{Split colliding labels.} Where the registry shows one label
with multiple registered semantics, require the qualified form
(``DAPO(dynamic-sampling)'') in abstracts and tables. Precedents exist: the
field already distinguishes ``PPO'' from ``PPO-clip'' when it matters
\citep{schulman2017proximal, henderson2018deep}.
\item \textbf{Reward the telemetry.} Per-step ZVF/GU trajectories and
group-size schedules should count as artifact contributions in their own
right, as evaluation-harness versioning now does
\citep{biderman2024lessons, pineau2020improving}.
\end{enumerate}
The deep-RL community needed a reporting reckoning \citep{henderson2018deep}, a
statistics paper \citep{colas2019hitchhiker, agarwal2021deep}, and a decade
to converge on reporting norms. RL-for-LLMs can shortcut that decade: the
confounds are already documented, the run manifest is seven fields, the
evaluation requirement is the eighth item, and enforcement is a diff. Report
the stack, not the label.
